Fix \(P\in\mathcal P_{d,\epsilon,M,\sigma}\), put \(L:=\log(en)\), and normalize
\[
 V:=Y/M,\qquad \nu_{ak}:=\mu_{ak}/M,\qquad
 s_{ak}:=P(X=k,A=a),\qquad
 z_{ak}:=E[V\mathbf1\{X=k,A=a\}],
\]
\[
 \varphi_k:=p_k(\nu_{1k}-\nu_{0k}).
\tag{1}
\]
Then \(\tau(P)/M=\sum_k\varphi_k\), \(z_{ak}=s_{ak}\nu_{ak}\), \(|\nu_{ak}|\le1/2\), \(s_{ak}\ge\epsilon p_k\), and
\[
 \operatorname{Var}(V\mid A=a,X=k)\le1.
\tag{2}
\]
Thus no outcome tail property beyond the declared conditional second moment is used.

For the calibrated branch, take exactly the split, degree \(K\), scale \(B\), pilot threshold, coefficients \(g_{Kj}\), light statistics \(\widehat P_k\), and heavy statistics \(\widehat h_k\) in Definition~\(\mathrm{def:polynomial\text{-}handle}\).  With continuous extensions at zero, define
\[
 H_K(x):=\frac{1-T_K(1-2x)}{2K^2},\qquad
 E_K(x):=\frac{H_K(x)}x,\qquad
 G_K(x):=\frac{1-E_K(x)}x.
\tag{3}
\]
Expanding the shifted Chebyshev polynomial gives \(G_K(x)=\sum_{j=0}^{K-2}g_{Kj}x^j\).  For \(0\le x\le1\), the identity \(T_K(\cos\theta)=\cos(K\theta)\) gives
\[
 0\le x\{1-xG_K(x)\}=H_K(x)\le K^{-2}.
\tag{4}
\]
If \(p_k\le B\), substitute \(x=s_{ak}/B\) and use \(s_{ak}\ge\epsilon p_k\) to obtain
\[
 \left|p_k\nu_{ak}-\frac{p_kz_{ak}}B G_K(s_{ak}/B)\right|
 \le\frac{B}{2\epsilon K^2}.
\tag{5}
\]
Independence of the ordered distinct indices yields
\[
 E[U_{k,a,j}]=z_{ak}p_ks_{ak}^j,
\tag{6}
\]
so \(\widehat P_k\) is unbiased for the difference of the two polynomial expressions in (5).

Let
\[
 B_-:=\frac{\lambda_0L}{2m_0},\qquad
 B_+:=\frac{2\lambda_0L}{m_0},
\]
and let \(\mathcal A\) be the event that every pilot-heavy cell has \(p_k\ge B_-\) and every pilot-light cell has \(p_k\le B_+\).  The balanced split and \(b_0=16\lambda_0\) give \(B_+\le B/4\).  The binomial Chernoff inequalities
\[
 P\{D\ge2tq\}\le(e/4)^{tq},\qquad
 P\{D\le tq/2\}\le e^{-tq/8}
\tag{7}
\]
and a union bound imply, uniformly for \(d\le\rho_{\epsilon}nL\),
\[
 Q_P^{(n)}(\mathcal A^c)\le2d e^{-c\lambda_0L}\le Cn^{-6}.
\tag{8}
\]
Increase \(N_{\epsilon}\), if needed, so that \(K\ge2\),
\[
 4(K+2)^2\le m_1,\qquad
 \frac{4(K+2)}{m_1}\le\frac{3B}{4}.
\tag{9}
\]
Conditional on the pilot and on \(\mathcal A\), the light set is deterministic relative to \(I_1\) and lies in \(\{k:p_k\le B/4\}\).  Equations~(5)--(6) and Lemma~\(\mathrm{lem:linear\text{-}mark\text{-}factorial\text{-}covariance}\) therefore give
\[
\begin{aligned}
 E\!\left[\left\{\sum_{k\in\widehat{\mathcal L}}\widehat P_k
 -\sum_{k\in\widehat{\mathcal L}}\varphi_k\right\}^2\middle|I_0\right]
 \le C_{\epsilon}\left[
 \frac1{m_1}+6^{2K}\left\{dB^2+\frac{d^2K^2B^2}{m_1}\right\}
 +\frac{d^2B^2}{K^4}\right].
\end{aligned}
\tag{10}
\]
The last term is the squared sum of the biases in (5).  Put \(\kappa:=2\alpha_0\log6\le1/32\).  Since \(6^{2K}\le n^{\kappa}\), \(B\le CL/n\), and \(K\le CL\), completion of squares gives
\[
 n^{\kappa}\frac{dL^2}{n^2}
 \le C\left\{\frac1n+\frac{d^2}{n^2L^2}\right\},
 \qquad
 n^{\kappa}\frac{d^2L^4}{n^3}
 \le C\frac{d^2}{n^2L^2},
\tag{11}
\]
and \(dB/K^2\le Cd/(nL)\).  Substitution into (10) proves
\[
 E\!\left[\left\{\sum_{k\in\widehat{\mathcal L}}\widehat P_k
 -\sum_{k\in\widehat{\mathcal L}}\varphi_k\right\}^2\middle|I_0\right]
 \le C_{\epsilon}\left\{\frac1n+\frac{d^2}{n^2L^2}\right\}.
\tag{12}
\]

For heavy cells, conditional on the pilot and all category-treatment indicators in \(I_1\), equation~(2) gives
\[
 \operatorname{Var}\!\left(\sum_{k\in\widehat{\mathcal H}}\widehat h_k
 \middle|(X_i,A_i)_{i\in I_1}\right)
 \le\frac1{m_1^2}\sum_{k\in\widehat{\mathcal H}}(N_k^{(1)})^2
 \left\{\frac{\mathbf1\{N_{1k}^{(1)}>0\}}{N_{1k}^{(1)}}+
 \frac{\mathbf1\{N_{0k}^{(1)}>0\}}{N_{0k}^{(1)}}\right\}.
\tag{13}
\]
For \(D\sim\operatorname{Binomial}(t,q)\),
\[
 E[(D+1)^{-1}]=\frac{1-(1-q)^{t+1}}{(t+1)q},
\tag{14}
\]
and \(D^{-1}\mathbf1\{D>0\}\le2(D+1)^{-1}\).  Overlap and conditioning on \(N_k^{(1)}\) show that the expectation of (13) is at most \(C_{\epsilon}/m_1\).

The conditional mean of the heavy sum equals
\[
 \frac1{m_1}\sum_{k\in\widehat{\mathcal H}}N_k^{(1)}
 (\nu_{1k}-\nu_{0k})-R_1+R_0,
\quad
 R_a:=\frac1{m_1}\sum_{k\in\widehat{\mathcal H}}
 N_k^{(1)}\nu_{ak}\mathbf1\{N_{ak}^{(1)}=0\}.
\tag{15}
\]
The first term is an empirical average bounded in absolute value by one, is unbiased for \(\sum_{k\in\widehat{\mathcal H}}\varphi_k\), and has variance at most \(1/m_1\).  Exact multinomial factorial moments, \(|\nu_{ak}|\le1/2\), and \(s_{ak}\ge\epsilon p_k\) give
\[
 E[R_a^2\mid I_0]\le C_{\epsilon}\left[
 \frac1{m_1}+
 \left\{\sum_{k\in\widehat{\mathcal H}}p_ke^{-c_{\epsilon}m_1p_k}\right\}^2
 \right].
\tag{16}
\]
On \(\mathcal A\), each selected heavy cell has \(p_k\ge B_-\ge cL/n\), so
\[
 E[R_a^2\mid I_0]\le C_{\epsilon}\left\{\frac1n+\frac{d^2}{n^2L^2}\right\}.
\tag{17}
\]
Equations~(13)--(17) give the heavy-cell analogue of (12).

The heavy and light sets partition the alphabet.  If \(Z\) is their normalized sum before clipping, then
\[
 E\!\left[\mathbf1_{\mathcal A}\left\{Z-\frac{\tau(P)}M\right\}^2\right]
 \le C_{\epsilon}\left\{\frac1n+\frac{d^2}{n^2L^2}\right\}.
\tag{18}
\]
Lemma~\(\mathrm{lem:scale\text{-}sanity}\) gives \(\tau(P)/M\in[-1,1]\), so projection cannot increase loss; the projected squared loss is at most four on \(\mathcal A^c\).  Equations~(8) and (18) prove the desired rate in the calibrated branch.

If \(d>\rho_{\epsilon}nL\), the declared zero fallback has loss at most \(M^2\), while
\[
 \min\left\{1,\frac{d^2}{n^2L^2}\right\}\ge\min(1,\rho_{\epsilon}^2).
\tag{19}
\]
If \(n<N_{\epsilon}\), the same conclusion follows from \(1\le N_{\epsilon}/n\).  Enlarging \(C_{\epsilon}\) proves the bound for every \(n,d\ge1\).  All divisions are indicator-totalized.  After aggregation each ordered factorial statistic is computed from cell counts, one marked sum, and falling factorials, giving \(O(dK^2)\) arithmetic operations in the calibrated branch.